\documentclass[handout]{beamer}
\usetheme{Madrid}
\usecolortheme{default}

\usepackage{amsmath}
\usepackage{amssymb}
\usepackage{tikz}
\usepackage[ngerman]{babel}

\title{Von klassischer zu moderner Kryptographie}
\subtitle{Einführung in die Public-Key-Kryptographie}
\author{Informatik}
\date{}

\begin{document}

\frame{\titlepage}

\begin{frame}{Die heutige Reise}
\tableofcontents
\end{frame}

\section{Rückblick: Klassische Chiffren}

\begin{frame}{Was wir bisher gelernt haben}
\begin{block}{Klassische Chiffren}
\begin{itemize}
    \item Caesar-Chiffre (Verschiebungschiffre)
    \item Vigenère-Chiffre (polyalphabetische Substitution)
    \item Transpositionschiffren
\end{itemize}
\end{block}

\pause

\begin{alertblock}{Frage}
Was haben all diese Chiffren gemeinsam?
\end{alertblock}
\end{frame}

\begin{frame}{Kritische Analyse: Caesar-Chiffre}
\begin{example}
Verschlüssele ``HALLO'' mit Verschiebung 7:
\[
\text{HALLO} \rightarrow \text{OHSSV}
\]
\end{example}

\pause

\begin{alertblock}{Schwachstelle}
Wenn ein Angreifer den Verschiebungswert kennt (oder alle 26 Möglichkeiten durchprobiert), ist die Chiffre vollständig gebrochen!
\end{alertblock}

\pause

\begin{block}{Wichtige Erkenntnis}
Die Sicherheit hängt davon ab, die \textbf{Methode} geheim zu halten, nicht nur den Schlüssel.
\end{block}
\end{frame}

\begin{frame}{Die Vigenère-Chiffre}
\begin{block}{Einst als ``unknackbar'' bezeichnet}
Die Vigenère-Chiffre galt jahrhundertelang als unknackbar.
\end{block}

\pause

\begin{alertblock}{Aber sie wurde geknackt!}
\begin{itemize}
    \item Häufigkeitsanalyse
    \item Kasiski-Untersuchung
    \item Schlüssellänge finden, dann als mehrere Caesar-Chiffren behandeln
\end{itemize}
\end{alertblock}

\pause

\begin{block}{Das Muster}
Alle klassischen Chiffren können mit genügend Geheimtext und Rechenaufwand gebrochen werden.
\end{block}
\end{frame}

\begin{frame}{Die zwei grundlegenden Probleme}
\begin{enumerate}
    \item \textbf{Sicherheit durch Geheimhaltung}
    \begin{itemize}
        \item Klassische Chiffren verlassen sich darauf, die Methode geheim zu halten
        \item Sobald die Methode bekannt ist, ist die Sicherheit gefährdet
    \end{itemize}

    \pause

    \item \textbf{Das Schlüsselverteilungsproblem}
    \begin{itemize}
        \item Wie können zwei Personen, die sich nie getroffen haben, einen geheimen Schlüssel austauschen?
        \item Wenn sie sich persönlich treffen, großartig! Aber was ist mit Fernkommunikation?
        \item Man kann den Schlüssel nicht über einen unsicheren Kanal senden...
    \end{itemize}
\end{enumerate}

\pause

\begin{alertblock}{Frage}
Wie können wir diese Probleme lösen?
\end{alertblock}
\end{frame}

\section{Kerckhoffs' Prinzip}

\begin{frame}{Auftritt: Kerckhoffs' Prinzip}
\begin{block}{Auguste Kerckhoffs (1883)}
\begin{quote}
\textit{``Ein Kryptosystem sollte sicher sein, auch wenn alles über das System, außer dem Schlüssel, öffentlich bekannt ist.''}
\end{quote}
\end{block}

\pause

\begin{alertblock}{Moment... Was?}
Wie kann ein System sicher sein, wenn jeder weiß, wie es funktioniert?
\end{alertblock}

\pause

\begin{block}{Die revolutionäre Idee}
Die Sicherheit sollte \textbf{nur} von der Geheimhaltung des Schlüssels abhängen, nicht vom Algorithmus!
\end{block}
\end{frame}

\begin{frame}{Warum Kerckhoffs' Prinzip wichtig ist}
\begin{block}{Vorteile öffentlicher Algorithmen}
\begin{itemize}
    \item Können von Experten weltweit untersucht und getestet werden
    \item Schwachstellen können gefunden und behoben werden
    \item Keine Notwendigkeit, Methoden geheim zu halten (in der Praxis unmöglich)
    \item Vertrauen durch Transparenz
\end{itemize}
\end{block}

\pause

\begin{exampleblock}{Moderne Realität}
Alle moderne Verschlüsselung (HTTPS, Banking, Messaging) verwendet öffentlich bekannte Algorithmen!
\end{exampleblock}

\pause

\begin{alertblock}{Aber wie?}
Die Antwort liegt in der \textbf{Mathematik}...
\end{alertblock}
\end{frame}

\section{Einwegfunktionen}

\begin{frame}{Die Grundlage: Einwegfunktionen}
\begin{definition}
Eine \textbf{Einwegfunktion} ist eine Funktion, die:
\begin{itemize}
    \item In einer Richtung leicht zu berechnen ist
    \item Extrem schwierig (praktisch unmöglich) umzukehren ist
\end{itemize}
\end{definition}

\pause

\begin{example}[Alltägliche Analogie]
\textbf{Farben mischen:}
\begin{itemize}
    \item \textcolor{green!60!black}{Einfach:} Gelb + Blau mischen = Grün
    \item \textcolor{red!60!black}{Schwer:} Bei Grün die genauen Gelb- und Blauanteile finden
\end{itemize}
\end{example}
\end{frame}

\begin{frame}{Mathematische Einwegfunktionen}
\begin{block}{Beispiel 1: Modulare Exponentiation}
\[
f(x) = g^x \bmod p
\]
wobei $g$ eine Basis und $p$ eine große Primzahl ist.
\end{block}

\pause

\begin{example}
Sei $g = 3$, $p = 17$, $x = 5$

\textbf{Vorwärts (Einfach):}
\[
3^5 \bmod 17 = 243 \bmod 17 = 5
\]

\pause

\textbf{Rückwärts (Schwer):}

Gegeben $g = 3$, $p = 17$, Ergebnis $= 5$, finde $x$ sodass $3^x \bmod 17 = 5$

Das ist das \textbf{diskrete Logarithmusproblem}!
\end{example}
\end{frame}

\begin{frame}{Ein weiteres Beispiel: Primfaktorzerlegung}
\begin{block}{Multiplikation vs. Faktorisierung}
\textbf{Einfache Richtung:}
\[
53 \times 79 = 4187
\]

\pause

\textbf{Schwere Richtung:}

Gegeben 4187, finde die zwei Primfaktoren.

\pause

Mit kleinen Zahlen: durch Ausprobieren möglich

Mit 400-stelligen Zahlen: würde Milliarden Jahre dauern!
\end{block}

\pause

\begin{exampleblock}{In der Praxis verwendet}
Diese Eigenschaft ist die Grundlage der RSA-Verschlüsselung!
\end{exampleblock}
\end{frame}

\begin{frame}{Warum Einwegfunktionen wichtig sind}
\begin{block}{Die zentrale Erkenntnis}
Wenn wir Kryptographie auf Einwegfunktionen aufbauen:
\begin{itemize}
    \item Kann der Algorithmus öffentlich sein (jeder kennt $g^x \bmod p$)
    \item Aber die Umkehrung ohne das Geheimnis ist unmöglich
    \item Sicherheit kommt von rechnerischer Schwierigkeit, nicht von Geheimhaltung
\end{itemize}
\end{block}

\pause

\begin{alertblock}{Das ändert alles}
Einwegfunktionen ermöglichen es uns, Systeme zu erstellen, die Kerckhoffs' Prinzip erfüllen!
\end{alertblock}
\end{frame}

\section{Diffie-Hellman-Schlüsselaustausch}

\begin{frame}{Der revolutionäre Durchbruch (1976)}
\begin{block}{Whitfield Diffie \& Martin Hellman}
Lösten das Schlüsselverteilungsproblem mit einer verblüffenden Idee:

\vspace{0.3cm}

\textit{Zwei Personen können ein gemeinsames Geheimnis über einen öffentlichen Kanal erstellen, ohne das Geheimnis selbst jemals zu übertragen!}
\end{block}

\pause

\begin{alertblock}{Denken Sie darüber nach}
Alice und Bob können sich auf einen geheimen Schlüssel einigen, während Eve alles mithört, was sie sagen!
\end{alertblock}
\end{frame}

\begin{frame}{Die Farbmisch-Analogie}
\begin{center}
\begin{tikzpicture}[scale=0.8]
    % Public color
    \draw[fill=yellow!80] (0,3) circle (0.4cm);
    \node at (0,2.3) {\small Öffentlich: Gelb};

    \pause

    % Alice's secret
    \draw[fill=red!80] (-3,1.5) circle (0.3cm);
    \node at (-3,0.9) {\small Alices Geheimnis: Rot};

    % Bob's secret
    \draw[fill=blue!80] (3,1.5) circle (0.3cm);
    \node at (3,0.9) {\small Bobs Geheimnis: Blau};

    \pause

    % Alice mixes
    \draw[fill=orange!80] (-3,0) circle (0.4cm);
    \node at (-3,-0.7) {\small Alice: Orange};
    \draw[->,thick] (-3,1.2) -- (-3,0.5);

    % Bob mixes
    \draw[fill=green!60!blue!80] (3,0) circle (0.4cm);
    \node at (3,-0.7) {\small Bob: Grün};
    \draw[->,thick] (3,1.2) -- (3,0.5);

    \pause

    % Exchange
    \draw[->,very thick,red] (-2.7,-0.3) to[bend left=20] (2.7,-0.3);
    \draw[->,very thick,blue] (2.7,0.3) to[bend left=20] (-2.7,0.3);
    \node at (0,-0.8) {\small \textbf{Öffentlicher Austausch}};

    \pause

    % Final mix
    \draw[fill=brown!70] (-3,-2) circle (0.4cm);
    \node at (-3,-2.7) {\small Alices Endfarbe};
    \draw[->,thick] (-3,-0.5) -- (-3,-1.5);

    \draw[fill=brown!70] (3,-2) circle (0.4cm);
    \node at (3,-2.7) {\small Bobs Endfarbe};
    \draw[->,thick] (3,-0.5) -- (3,-1.5);

    \pause

    \node at (0,-3.2) {\textbf{Gleiche Farbe!}};
\end{tikzpicture}
\end{center}
\end{frame}

\begin{frame}{Das mathematische Protokoll}
\begin{block}{Öffentliche Parameter (allen bekannt)}
\begin{itemize}
    \item $p$ = eine große Primzahl
    \item $g$ = ein Generator (eine Zahl zwischen 1 und $p-1$)
\end{itemize}
\end{block}

\pause

\begin{block}{Das Protokoll}
\begin{enumerate}
    \item Alice wählt Geheimnis $a$, berechnet $A = g^a \bmod p$, sendet $A$ an Bob
    \item Bob wählt Geheimnis $b$, berechnet $B = g^b \bmod p$, sendet $B$ an Alice
    \item Alice berechnet: $s = B^a \bmod p = (g^b)^a \bmod p = g^{ab} \bmod p$
    \item Bob berechnet: $s = A^b \bmod p = (g^a)^b \bmod p = g^{ab} \bmod p$
\end{enumerate}
\end{block}

\pause

\begin{exampleblock}{Ergebnis}
Beide haben das gleiche gemeinsame Geheimnis $s = g^{ab} \bmod p$!
\end{exampleblock}
\end{frame}

\begin{frame}{Warum ist das sicher?}
\begin{block}{Was Eve (die Lauscherin) weiß:}
\begin{itemize}
    \item $p$ (öffentliche Primzahl)
    \item $g$ (öffentlicher Generator)
    \item $A = g^a \bmod p$ (abgefangen)
    \item $B = g^b \bmod p$ (abgefangen)
\end{itemize}
\end{block}

\pause

\begin{alertblock}{Was Eve braucht, um das Geheimnis zu finden:}
Sie muss lösen: $g^a \bmod p = A$ um $a$ zu finden (oder das Äquivalent für $b$)

Das ist das \textbf{diskrete Logarithmusproblem}!
\end{alertblock}

\pause

\begin{block}{Mit großen Zahlen}
Mit 2048-Bit-Primzahlen würde dies länger als das Alter des Universums dauern!
\end{block}
\end{frame}

\section{Durchgerechnetes Beispiel}

\begin{frame}{Vollständiges Beispiel: Einrichtung}
\begin{block}{Öffentliche Parameter}
\begin{align*}
p &= 23 \quad \text{(Primzahl)} \\
g &= 5 \quad \text{(Generator)}
\end{align*}
\end{block}

\pause

\begin{block}{Alices Zug}
Alice wählt Geheimnis: $a = 6$

Alice berechnet:
\[
A = g^a \bmod p = 5^6 \bmod 23
\]

\pause

\[
A = 15625 \bmod 23 = 8
\]

\textbf{Alice sendet $A = 8$ an Bob (öffentlich)}
\end{block}
\end{frame}

\begin{frame}{Vollständiges Beispiel: Bobs Zug}
\begin{block}{Bobs Zug}
Bob wählt Geheimnis: $b = 15$

Bob berechnet:
\[
B = g^b \bmod p = 5^{15} \bmod 23
\]

\pause

\[
B = 30517578125 \bmod 23 = 19
\]

\textbf{Bob sendet $B = 19$ an Alice (öffentlich)}
\end{block}

\pause

\begin{alertblock}{Bisheriger Stand}
Öffentlich bekannt: $p = 23$, $g = 5$, $A = 8$, $B = 19$

Geheimnisse: Alice kennt $a = 6$, Bob kennt $b = 15$
\end{alertblock}
\end{frame}

\begin{frame}{Vollständiges Beispiel: Berechnung des gemeinsamen Geheimnisses}
\begin{block}{Alice berechnet}
\[
s_A = B^a \bmod p = 19^6 \bmod 23
\]

\pause

\[
s_A = 47045881 \bmod 23 = 2
\]
\end{block}

\pause

\begin{block}{Bob berechnet}
\[
s_B = A^b \bmod p = 8^{15} \bmod 23
\]

\pause

\[
s_B = 35184372088832 \bmod 23 = 2
\]
\end{block}

\pause

\begin{exampleblock}{Erfolg!}
Sowohl Alice als auch Bob haben das gleiche gemeinsame Geheimnis berechnet: $s = 2$
\end{exampleblock}
\end{frame}

\begin{frame}{Was ist mit Eve?}
\begin{block}{Eves Informationen}
\begin{itemize}
    \item $p = 23$
    \item $g = 5$
    \item $A = 8$
    \item $B = 19$
\end{itemize}
\end{block}

\pause

\begin{alertblock}{Eves Herausforderung}
Um das Geheimnis $s = g^{ab} \bmod 23 = 2$ zu finden, muss sie lösen:
\[
5^a \bmod 23 = 8 \quad \text{oder} \quad 5^b \bmod 23 = 19
\]
\end{alertblock}

\pause

\begin{block}{Realitätscheck}
\begin{itemize}
    \item Mit $p = 23$: Eve kann alle 22 Möglichkeiten durchprobieren (einfach)
    \item Mit 2048-Bit $p$: Eve würde Milliarden Jahre brauchen!
\end{itemize}
\end{block}
\end{frame}

\begin{frame}{Verifikation: Schritt für Schritt}
\begin{block}{Überprüfen wir Alices Berechnung}
\begin{align*}
s_A &= B^a \bmod p \\
    &= 19^6 \bmod 23 \\
    &= (19^2)^3 \bmod 23 \\
    &= 361^3 \bmod 23 \\
    &= (361 \bmod 23)^3 \bmod 23 \\
    &= 16^3 \bmod 23 \\
    &= 4096 \bmod 23 \\
    &= 2
\end{align*}
\end{block}
\end{frame}

\begin{frame}{Die Magie enthüllt}
\begin{block}{Warum es funktioniert}
\begin{align*}
s_A &= B^a \bmod p = (g^b)^a \bmod p = g^{ab} \bmod p \\
s_B &= A^b \bmod p = (g^a)^b \bmod p = g^{ab} \bmod p
\end{align*}

Beide berechnen $g^{ab} \bmod p$ auf unterschiedlichen Wegen!
\end{block}

\pause

\begin{exampleblock}{Der Durchbruch}
\begin{itemize}
    \item Alice und Bob haben das Geheimnis nie ausgetauscht
    \item Jeder nutzte seine privaten Informationen mit dem öffentlichen Wert des anderen
    \item Eve kann $g^{ab}$ nicht berechnen, ohne $a$ oder $b$ zu kennen
\end{itemize}
\end{exampleblock}
\end{frame}

\begin{frame}{Übungsaufgabe}
\begin{block}{Sie sind dran!}
Gegeben: $p = 11$, $g = 2$

Alices Geheimnis: $a = 3$

Bobs Geheimnis: $b = 7$

\vspace{0.5cm}

Berechnen Sie:
\begin{enumerate}
    \item $A$ (Alices öffentlicher Wert)
    \item $B$ (Bobs öffentlicher Wert)
    \item $s$ (das gemeinsame Geheimnis)
\end{enumerate}
\end{block}

\pause

\begin{alertblock}{Hinweis}
Denken Sie daran, bei jedem Schritt den Modulo zu nehmen, um die Zahlen handhabbar zu halten!
\end{alertblock}
\end{frame}

\begin{frame}{Übungsaufgabe: Lösung}
\begin{block}{Berechnungen}
\begin{align*}
A &= g^a \bmod p = 2^3 \bmod 11 = 8 \bmod 11 = 8 \\
B &= g^b \bmod p = 2^7 \bmod 11 = 128 \bmod 11 = 7 \\
\end{align*}

\pause

\begin{align*}
s_{\text{Alice}} &= B^a \bmod p = 7^3 \bmod 11 = 343 \bmod 11 = 2 \\
s_{\text{Bob}} &= A^b \bmod p = 8^7 \bmod 11 = 2097152 \bmod 11 = 2
\end{align*}
\end{block}

\pause

\begin{exampleblock}{Antwort}
Gemeinsames Geheimnis: $s = 2$
\end{exampleblock}
\end{frame}

\section{Zusammenfassung}

\begin{frame}{Was wir heute gelernt haben}
\begin{enumerate}
    \item \textbf{Klassische Chiffren} haben grundlegende Schwächen
    \begin{itemize}
        \item Sicherheit durch Geheimhaltung
        \item Schlüsselverteilungsproblem
    \end{itemize}

    \pause

    \item \textbf{Kerckhoffs' Prinzip}: Sicherheit sollte nur vom Schlüssel abhängen

    \pause

    \item \textbf{Einwegfunktionen}: Vorwärts einfach, rückwärts schwer
    \begin{itemize}
        \item Modulare Exponentiation
        \item Primfaktorzerlegung
    \end{itemize}

    \pause

    \item \textbf{Diffie-Hellman}: Zwei Parteien können ein gemeinsames Geheimnis über einen öffentlichen Kanal erstellen!
\end{enumerate}
\end{frame}

\begin{frame}{Das große Ganze}
\begin{block}{Von klassisch zu modern}
\begin{itemize}
    \item \textbf{Klassisch}: Die Methode geheim halten
    \item \textbf{Modern}: Nur den Schlüssel geheim halten
\end{itemize}
\end{block}

\pause

\begin{block}{Die Revolution}
Mathematische Einwegfunktionen ermöglichen es uns:
\begin{itemize}
    \item Öffentliche Algorithmen zu verwenden
    \item Das Schlüsselverteilungsproblem zu lösen
    \item Sicherheit durch Berechnung statt Geheimhaltung zu erreichen
\end{itemize}
\end{block}

\pause

\begin{exampleblock}{Auswirkungen in der realen Welt}
Jedes Mal, wenn Sie ``https://'' in Ihrem Browser sehen, nutzen Sie diese Prinzipien!
\end{exampleblock}
\end{frame}

\begin{frame}{Ausblick}
\begin{block}{Was kommt als Nächstes?}
\begin{itemize}
    \item RSA-Verschlüsselung (ein weiteres Public-Key-System)
    \item Digitale Signaturen
    \item Hybride Kryptosysteme
    \item Bedrohungen durch Quantencomputer
\end{itemize}
\end{block}

\pause

\begin{alertblock}{Denken Sie daran}
Moderne Kryptographie basiert auf soliden mathematischen Grundlagen, die Sicherheit durch Transparenz ermöglichen!
\end{alertblock}
\end{frame}

\begin{frame}
\begin{center}
{\Huge Fragen?}

\vspace{1cm}

\large
Denken Sie daran: Das gemeinsame Geheimnis wird nie übertragen!
\end{center}
\end{frame}

\end{document}
